\section{Conclusion}

Tool-augmented LLM agents expose a security surface that extends well beyond user-input filtering. In practical gateway deployments, risk emerges across message ingress, prompt construction, tool invocation, tool-returned content, outbound messaging, and local control files. These distributed attack surfaces render single-boundary defenses insufficient for many real agent workflows, particularly when prompt injection, tool abuse, credential leakage, and stateful escalation can interact across multiple runtime stages.

This paper has presented OpenClaw PRISM, a zero-fork runtime security layer for OpenClaw-based agent gateways. Rather than proposing a new detection model, PRISM combines ten lifecycle hooks, a hybrid heuristic-plus-LLM scanning pipeline, conversation-scoped and session-scoped risk accumulation, policy-enforced controls over tools, paths, and outbound traffic, and a tamper-evident audit and operations plane. The resulting architecture aims to advance agent security from isolated input filtering toward a deployable runtime layer that can be attached to an existing gateway, supervised as a set of long-running services, audited after the fact, and adjusted through explicit operational workflows---all without forking the upstream framework.

The broader message of this work is that agent security should be treated as a systems problem as much as a detection problem. Runtime interception, staged policy response, audit integrity, and recoverable operator workflows all matter when security controls are expected to function inside long-running, tool-using gateways. PRISM is one concrete instantiation of that perspective: not a complete defense, not a formal guarantee, and not a substitute for lower-layer hardening such as OS sandboxing or network-level egress control, but a practical architecture for making agent gateways more controllable, observable, and defensible in deployment.

Several directions remain open for future work. Stronger benchmark coverage---especially for indirect injection, tool abuse, and long-horizon escalation---would improve the evidence base for runtime defenses of this kind. Extending the architecture beyond the OpenClaw ecosystem would help clarify which mechanisms are framework-specific and which generalize to other agent runtimes. A particularly promising direction is to pair PRISM's deployable hook and policy infrastructure with more formal safe-tool-use specifications or with stronger audit-driven policy refinement. Together, these directions point toward a broader agenda in which agent security is constructed from layered runtime controls, explicit policy boundaries, and operationally verifiable system behavior.
